%%%%%%%%%%%%%%%%%%%%%%%%%%%%%%%%%%%%%%%%%
% Programming/Coding Assignment
% LaTeX Template
%
% This template has been downloaded from:
% http://www.latextemplates.com
%
% Original author:
% Ted Pavlic (http://www.tedpavlic.com)
%
% Note:
% The \lipsum[#] commands throughout this template generate dummy text
% to fill the template out. These commands should all be removed when 
% writing assignment content.
%
% This template uses a Perl script as an example snippet of code, most other
% languages are also usable. Configure them in the "CODE INCLUSION 
% CONFIGURATION" section.
%
%%%%%%%%%%%%%%%%%%%%%%%%%%%%%%%%%%%%%%%%%

%----------------------------------------------------------------------------------------
%	PACKAGES AND OTHER DOCUMENT CONFIGURATIONS
%----------------------------------------------------------------------------------------

\documentclass{article}

\usepackage{fancyhdr} % Required for custom headers
\usepackage{lastpage} % Required to determine the last page for the footer
\usepackage{extramarks} % Required for headers and footers
\usepackage[usenames,dvipsnames]{color} % Required for custom colors
\usepackage{graphicx} % Required to insert images
\usepackage{listings} % Required for insertion of code
\usepackage{courier} % Required for the courier font
\usepackage{multirow}
\usepackage{hyperref}


% Margins
\topmargin=-0.45in
\evensidemargin=0in
\oddsidemargin=0in
\textwidth=6.5in
\textheight=9.0in
\headsep=0.25in

\linespread{1.1} % Line spacing

%----------------------------------------------------------------------------------------
%	CODE INCLUSION CONFIGURATION
%----------------------------------------------------------------------------------------

\definecolor{MyDarkGreen}{rgb}{0.0,0.4,0.0} % This is the color used for comments
\lstloadlanguages{c} % Load Perl syntax for listings, for a list of other languages supported see: ftp://ftp.tex.ac.uk/tex-archive/macros/latex/contrib/listings/listings.pdf
\lstset{language=[sharp]c, % Use Perl in this example
        frame=single, % Single frame around code
        basicstyle=\small\ttfamily, % Use small true type font
        keywordstyle=[1]\color{Blue}\bf, % Perl functions bold and blue
        keywordstyle=[2]\color{Purple}, % Perl function arguments purple
        keywordstyle=[3]\color{Blue}\underbar, % Custom functions underlined and blue
        identifierstyle=, % Nothing special about identifiers                                         
        commentstyle=\usefont{T1}{pcr}{m}{sl}\color{MyDarkGreen}\small, % Comments small dark green courier font
        stringstyle=\color{Purple}, % Strings are purple
        showstringspaces=false, % Don't put marks in string spaces
        tabsize=5, % 5 spaces per tab
        %
        % Put standard Perl functions not included in the default language here
        morekeywords={rand},
        %
        % Put Perl function parameters here
        morekeywords=[2]{on, off, interp},
        %
        % Put user defined functions here
        morekeywords=[3]{test},
       	%
        morecomment=[l][\color{Blue}]{...}, % Line continuation (...) like blue comment
        numbers=left, % Line numbers on left
        firstnumber=1, % Line numbers start with line 1
        numberstyle=\tiny\color{Blue}, % Line numbers are blue and small
        stepnumber=5 % Line numbers go in steps of 5
}

\newcommand{\horrule}[1]{\rule{\linewidth}{#1}}

% Creates a new command to include a perl script, the first parameter is the filename of the script (without .pl), the second parameter is the caption
\newcommand{\perlscript}[2]{
\begin{itemize}
\item[]\lstinputlisting[caption=#2,label=#1]{#1.cs}
\end{itemize}
}

\begin{document}

\begin{tabular}{l l}
\multirow{5}{*}{\includegraphics[width=2cm]{../../recursos/logo.png}} & Universidad del Istmo de Guatemala \\
 & Facultad de Ingenieria \\
 & Ing. en Sistemas \\
 & Informatica 1 \\
 & Prof. Ernesto Rodriguez - \href{mailto:erodriguez@unis.edu.gt}{erodriguez@unis.edu.gt} \\
\end{tabular}
\\\\\\

\begin{center}
        \horrule{0.5pt}
        \huge{Hoja de trabajo \#1} \\
        \large{Fecha de entrega: 25 de Julio, 2018 - 11:59pm} \\
        \horrule{1pt}
\end{center}

\emph{Instrucciones: Resolver cada uno de los ejercicios siguiendo sus respectivas
instrucciones. El trabajo debe ser entregado a traves de Github, en su repositorio del curso, colocado en una carpeta llamada "Hoja de trabajo 1".
Al menos que la pregunta indique diferente, todas las respuestas a preguntas escritas deben presentarse en
un documento formato pdf, el cual haya sido generado mediante Latex. }

% \perlscript{homework_example}{Sample Perl Script With Highlighting}

\section*{Iniciaci\'on}
\begin{enumerate}
        \item{Si no la tiene, crear una cuenta en \href{https://github.com/}{Github}.}
        \item{Crear un repositorio\cite{GitRepo} para los deberes correspondientes a este curso en Github.}
        \item{Enviar su usuario de Github y link de su repositorio al profesor
        \href{mailto:erodriguez@unis.edu.gt}{erodriguez@unis.edu.gt}.}
        \item{Tomar en cuenta que estea hoja
        (y las hojas restantes del curos) deben ser entregadas mediante un commit
        en el repositorio creado en el articulo 2 de Github. En la fecha y hora de entrega,
        el profesor copiara el estado acutal del repositorio y calificara ese material. Por
        favor asegurarse que todos los cambios esten \emph{empujados} (\texttt{git push}) en el
        repositorio antes de la fecha de entrega. Esto se puede revisar accediendo el sitio
        Github y ver que todos los archivos necesarios esten presentes en el repositorio.}
        \item{El ``Ejercicio \#1'' y ``Ejercicio \#4'' le ayudaran a publicar correctamente
        el trabajo en Github. Si aun tiene dudas, por favor contactar al profesor.}
        \item{Para hacer esta hoja de trabajo, necesitara Latex\cite{LatexPage}. Para trabajar
        con archivos Latex (``*.tex'') se recomienda utilizar un editor. La \href{https://marketplace.visualstudio.com/items?itemName=James-Yu.latex-workshop}{extension de
        Latex} para \emph{Visual Studio Code} es excelente, pero t\'ambien se puede utilizar
        un programa como \href{http://www.xm1math.net/texmaker/}{Texmaker}. Asegurese
        de configurar Latex de manera que pueda producir archivos \emph{PDF}.}
        \item{El Wiki de Latex\cite{Latex} es un buen lugar para consultar el formato.}
\end{enumerate}

\section*{Ejercicio \#1 (5\%)}
\begin{enumerate}
        \item{
                Si no lo ha hecho, crear un repositorio en Github para alojar
                las hojas de trabajo.}
        \item{
                Luego de haber creado un repositorio, debe clonar el repositorio
                a una carpeta local mediante el comando \texttt{git clone [url]} el cual
                puede ser ejecutado desde la terminal (Linux y Mac) o desde la consola
                \emph{Git Bash} (en Windows). Para mayor
                facilidad, utilizar un url tipo \emph{https}. Sin embargo para mayor
                seguridad y conveniencia se recomienda usar \emph{ssh}\cite{GithubSsh}.
                Como ejemplo, este comando clonaria el repositorio de este curso: \\
                \texttt{git clone https://github.com/netogallo/informatica-2018-2019}
        }
        \item{
                Luego de haber clonado el repositorio, se creara una carpeta
                con el nombre del repositorio dentro de la carpeta donde se
                haya ejecutado el comando anterior.
        }
\end{enumerate}

\section*{Ejercicio \#2: Abstracci\'on (50\%)}
Consideremos un dado de seis caras colocado en una superficie:
\begin{itemize}
        \item{La cara opuesta a la superficie muestra el \emph{numero actual} en que se encuentra el dado.}
        \item{El dado puede cambiar su \emph{numero actual} rotando 90 grados en cualquiera de
        sus 4 lados. La imagen muestra el dado cambiando su \emph{nuero actual} de ``2'' a ``6''
        rotando sobre la cara ``1''.}
\end{itemize}
\begin{center}
\includegraphics[width=3cm]{die.png}
\end{center}
Su deber es modelar el dado y sus transiciones de estado como un \emph{grafo}. Para ello debe
indicar:
\begin{enumerate}
        \item{El \emph{conjunto} de nodos del grafo.}
        \item{El \emph{conjunto} de vertices del grafo.}
\end{enumerate}
{\bf Pista:} Por lo general, los vertices corresponden a transiciones validas que pueden
ocurrir en el objeto que esta siendo modelado.

\section*{Ejercicio \#3 (40\%)}
Indique como puede utilizar el \emph{grafo} creado en el paso anterior para simular
un lanzamiento de dados. Supongamos, que el dado siempre inicia en el \emph{numero actual}
``1'':
\begin{itemize}
        \item{¿Que estructura de datos podria representar un lanzamiento de dados?}
        \item{¿Que algoritmo podriamos utilizar para generar dicha estructura?}
        \item{¿Como nos aseguramos que ese algoritmo siempre produce un resultado?}
\end{itemize}

\section*{Ejercicio \#4 (5\%)}
Luego de haber terminado los ejercicios anteriores, su trabajo debe
ser publicado a github. Para publicar a Github utilizando Git se procede
de la siguiente manera:
\begin{enumerate}
        \item {
                El comando \texttt{git status} listara todos los
                cambios pendientes por agregar al repositorio. Si no
                ha agregado ningun archivo al repositorio, listara todas
                las carpetas y el archivo \emph{.gitignore} dentro de la
                carpeta del proyecto.
        }
        \item{
                El comando \texttt{git add [archivos]} permite agregar
                archivos a la siguiente revision de git. Como atajo,
                puede ejecutar \texttt{git add .} para agregar todos
                los cambios que hay en el repositorio.
        }
        \item {
                El comando \texttt{git status} ahora indicara los archivos
                nuevos que hay en el repositorio. Si su archivo \emph{.gitignore}
                esta formateado correctamente, solo deberian agregarse:
                \begin{itemize}
                        \item{El archivo \emph{.gitignore}}
                        \item{El archivo \emph{pdf} correspondiente a esta hoja de trabajo}
                        \item{El archivo \emph{tex} correspondiente a esta hoja de trabajo}
                        \item{Un posible archivo o carpeta de configuraci\'on del editor.
                        El caso de ``Visual Studio Code'' sera una carpeta llamada \emph{.vscode}.}
                \end{itemize}
        }
        \item{
                Si el paso anterior muesta archivos extra (``*.aux'', ``*.log'', ``*.out'', ect...),
                debera componer su archivo \emph{.gitignore}. Para hacerlo, primero
                debe ejecutar el comando \texttt{git reset HEAD .}, el cual quitara
                todos los archivos de la lista de cambios. Luego modifique su archivo \emph{.gitignore}
                adecuadamente y vuelba a ejecutar \texttt{git add .}. Repetir hasta
                que la lista de archivos nuevos corresponda a la indicada en el paso \#3.
        }
\end{enumerate}


\bibliography{../../recursos/referencias}
\bibliographystyle{plain}

\end{document}